\documentclass[10pt,a4paper]{amsart}
\usepackage[margin=19mm]{geometry}
\usepackage{amsmath,amssymb,mathtools}
\usepackage[scr=rsfs,cal=euler]{mathalfa}
\usepackage{xfrac}
\usepackage[unicode,hidelinks,bookmarks=false]{hyperref}
\newcommand{\ie}{i.e.\ }
\newcommand{\cf}{cf.\ }
\newcommand{\resp}{respectively\ }
\newcommand{\Subsection}{\subsection}
\providecommand{\nobreakdash}{\nobreak\hbox{-}}
\setlength{\emergencystretch}{2em}
\multlinegap0pt
\newtheorem{corollary}{Corollary}
\newtheorem{lemma}{Lemma}
\newtheorem{theorem}{Theorem}
\newcommand{\abs}[1]{\left\vert #1 \right\vert}
\newcommand{\boldA}{\mathbf{A}}
\newcommand{\norm}[1]{\left\| #1 \right\|}
\title{The computable functional calculus}
\author{Christopher J. Eagle, Timothy H. McNicholl}
\date{Source: arXiv:2606.05456v1 (CC BY 4.0)}
\begin{document}
\maketitle

\noindent\emph{Source and licence.} The computable functional calculus. Version arXiv:2606.05456v1 (CC BY 4.0), distributed by arXiv under the Creative Commons Attribution 4.0 International licence, http://creativecommons.org/licenses/by/4.0/, as recorded in the arXiv licence metadata for this exact version and shown on the abstract page https://arxiv.org/abs/2606.05456v1. Full licence notice: \texttt{licenses/arxiv-2606.05456-CC-BY-4.0.txt}.\par\smallskip

\noindent\textit{Editorial scope.} Counted unit: the article's theorem thm:ApproximateUnit (source0.tex lines 365--367). The article's complete original proof of it is delivered with it (source0.tex lines 368--396, one \texttt{proof} environment of 29 lines). No mathematics is written, repaired or re-derived: the statement and the proof are the article's own lines, in the article's own notation, and no retained line is substituted or re-flowed.  Retained uncounted as the source-local context the proof needs: lines 307--309; lines 338--340; lines 347--349.  Nothing was omitted from any retained span, so the package carries no omission record.  No retained line contains a citation, so the wrapper carries no bibliography and no bibliography entry.  No retained line contains a figure environment, an image-inclusion command or drawing code, so no image is delivered and the package declares no asset dependency.  Editorial formatting only, in addition: an AMS \texttt{amsart} preamble that restates the article's own declarations and macros used by the retained lines, namely the environments \texttt{corollary}, \texttt{lemma}, \texttt{theorem} on separate counters, together with the standard packages those lines need.  The retained environments print in this document as Corollary 1, Lemma 1, Theorem 1 in order of appearance, whereas the article prints the same items with the numbering of its own full document.  The complete native source of the article as distributed in the arXiv e-print for this exact version is bundled unchanged as \texttt{sources/202118/source0.tex}.
\begin{corollary}\label{corollary:PositiveNegativeParts}
If $a \in \boldA$ is a normal computable point of $\boldA^\#$ then the real and imaginary parts $a_r$ and $a_i$ of $a$ are computable points of $\boldA^\#$.  If $a$ is a self-adjoint computable point of $\boldA^\#$ then the positive and negative parts $a^+$, and $a^-$ are computable points of $\boldA^\#$.
\end{corollary}

\begin{lemma}\label{lemma:ComputableApproxUnitDirected}
If $a, b \in \Lambda$ are computable points of $\boldA^\#$ then $\Lambda$ contains a computable point $c$ of $\boldA^\#$ such that $a \leq c$ and $b \leq c$.  Moreover, an $\boldA^\#$-index of $c$ can be computed from $\boldA^\#$-indices of $a$ and $b$.
\end{lemma}

\begin{lemma}\label{lemma:ComputableApproxUnitOneStep}
Suppose that $a \in \Lambda$, $a$ is a computable point of $\boldA^\#$, and $n \in \mathbb{N}$.  Then $\Lambda$ contains a computable point $e$ of $\boldA^\#$ such that whenever $c \in \Lambda$ and $c \geq e$ we have $\max\{\norm{a-ca}, \norm{a-ac}\} < 2^{-n}$.  Moreover, an $\boldA^\#$-index of $e$ can be computed from $n$ and an $\boldA^\#$-index of $a$.
\end{lemma}

\begin{theorem}\label{thm:ApproximateUnit}
There is an effective approximate unit of $\boldA^\#$.
\end{theorem}

\begin{proof}
Fix a c.e. set $\Lambda_0$ as in the previous lemma.  We recursively define a sequence $(u_n)_{n\in\mathbb{N}}$.  Suppose that $u_j$ has been defined for all $j < n$.  By Lemma \ref{lemma:ComputableApproxUnitOneStep}, for each $k \in \{0, \ldots, n\}$, we can compute an $\boldA^\#$-index of an $e_{k,n} \in \Lambda_0$ so that $\max\{\norm{c-ce_{k,n}}, \norm{c-e_{k,n}c}\} < 2^{-n}$ whenever $c \in \Lambda$ is such that $c \geq e_{k,n}$.  By Lemma \ref{lemma:ComputableApproxUnitDirected} we can then compute an $\boldA^\#$-index of a $u_n \in \Lambda$ such that $u_n \geq e_{k, n}$ for all $k \leq n$ and so that $u_n \geq u_j$ for all $j < n$.  We will show that this sequence $(u_n)_{n\in\mathbb{N}}$ is an effective approximate unit.

Fix a generated point $a$ of $\boldA^\#$.  Let $b = (\norm{a}+1)^{-1}a$.  Let $b_r$ and $b_i$ be the real and imaginary parts of $b$, respectively, and let their respective positive and negative parts be $b_r^+, b_r^-, b_i^+$, and $b_i^-$.  By Corollary \ref{corollary:PositiveNegativeParts}, $b_r^+, b_r^-, b_i^+$, and $b_i^-$ are computable points of $\boldA^\#$.  Moreover, it is possible to compute an $\boldA^\#$-index for each of them from a $\boldA^\#$-index for $a$.

Fix $k \in \mathbb{N}$.  Set $k_0 = -\lceil \log_2(\norm{a} + 1) \rceil + k + 4$.  Compute $k(r, +) \in \mathbb{N}$ such that $\norm{b_r^+-u_{k(r,+)}} < 2^{-k_0}$, and likewise compute such $k(r, -)$, $k(i, +)$, and $k(i, -)$.  Let $N_0 = \max\{k(r,+), k(r,-), k(i,+), k(i, -)\}$.

Fix $n \geq N_0$.  Since $\norm{u_n} < 1$ and $N_0 \geq k(r,+)$,
\begin{align*}
\abs{\norm{b_r^+ - u_n b_r^+} - \norm{u_{k(r,+)} - u_n u_{k(r,+)}}}
&= \abs{\norm{b_r^+} - \norm{u_{k(r,+)}}}\norm{1_{\boldA_1}-u_n} \\ &\leq \norm{b_r^+-u_{k(r,+)}}(\norm{1_{\boldA_1}}+\norm{u_n}) \\
&< 2\norm{b_r^+ - u_{k(r,+)} } \\
&< 2\cdot2^{-k_0}
\end{align*}
thus
\[\norm{b_r^+ - u_nb_r^+} < \norm{u_{k(r,+)} - u_nu_{k(r,+)}} + 2\cdot2^{-k_0}.\]
Also, by construction, $\norm{u_{k(r,+)}-u_nu_{k(r,+)}} < 2^{-k_0}$, so we obtain
\[\norm{b_r^+ - u_nb_r^+} < 3\cdot2^{-k_0}.\]
Similar calculations produce the same estimate using $b_r^-$, $b_i^+$, and $b_i^-$. The decomposition $b = b_r^+ - b_r^- + ib_i^+ - ib_i^-$
therefore implies
\[
\norm{b-u_nb} < 4\cdot3\cdot2^{-k_0} < 2^{-k_0+4}.
\]
Thus, by definition of $k_0$,
\[\norm{b-u_nb} < (\norm{a}+1)2^{-k}.\]
Finally, since $b = (\norm{a}+1)^{-1}a$, this implies the desired
\[\norm{a-u_na} < 2^{-k}.\]
It follows similarly that $\norm{a - a u_n} < 2^{-k}$.
\end{proof}

\end{document}
